\begin{abstract}
Groundbreaking advancements, exemplified by BitNet~\cite{bitnet}, have ushered in an era of 1-bit Large Language Models (LLMs). In this study, we present a novel 1-bit LLM architecture, designated as \textbf{\text{BitNet b1.58}}, which constrains every model parameter to one of three values: \{-1, 0, 1\}. Remarkably, this approach achieves parity with standard full-precision (FP16 or BF16) Transformer LLMs of equivalent size and trained with the same token count, in terms of perplexity and downstream task effectiveness. Simultaneously, it yields substantial benefits in latency, memory usage, throughput, and energy efficiency. Importantly, the 1.58-bit LLM establishes an alternative scaling law and provides new training methodologies for developing the next wave of high-efficiency, high-performance LLMs. Additionally, this technology suggests new computational strategies and paves the way for hardware tailored specifically for 1-bit LLM operations.
\end{abstract}

\section{The Era of 1-bit LLMs}

AI research has witnessed unprecedented progress in developing ever larger and more capable Large Language Models (LLMs). Despite their impressive results across a spectrum of natural language processing applications, the surging parameter counts have brought forth considerable obstacles for deployment, alongside rising concerns regarding energy consumption and sustainability.

A promising solution is post-training quantization, where model parameters and activations are compressed to lower bit-widths to streamline inference~\cite{smoothquant, gptq, quip, quip_sharp}. By minimizing numerical precision, these techniques dramatically lower memory footprints and computational load for LLMs. The industry has seen an evolution from 16-bit models toward even smaller bit-widths, such as the 4-bit options~\cite{gptq, awq}. Nonetheless, while post-training quantization remains prevalent in deployed LLMs, it is not the most optimal approach.

Recent advancements in 1-bit neural architectures, exemplified by BitNet~\cite{bitnet}, demonstrate a viable approach for decreasing the operational cost of LLMs without sacrificing accuracy. Standard LLMs utilize 16-bit floating-point representations (such as FP16 or BF16), and their computational workload is dominated by matrix multiplications. As a result, the primary computational expenses stem from floating-point operations—specifically, addition and multiplication. By contrast, BitNet’s matrix multiplications are executed solely via integer additions, leading to substantial energy savings when deploying LLMs. Given that power constraints often set the performance ceiling for hardware, reductions in energy consumption directly facilitate accelerated computation.

Beyond raw computation, transferring model weights from DRAM to on-chip memory (such as SRAM) imposes a considerable burden during inference. Some efforts have sought to increase SRAM capacity to enhance throughput, but the corresponding cost far exceeds that of DRAM. When contrasted with their full-precision counterparts, 1-bit LLMs feature significantly smaller memory footprints with respect to both storage and bandwidth, which dramatically lessens the resources and time required for weight transfer from DRAM, culminating in swifter and more economical inference procedures.

In this study, we present a new 1-bit LLM extension termed \textbf{\text{BitNet b1.58}}, where each parameter is ternary and may assume values from \{-1, 0, 1\}. By incorporating the 0 value into the 1-bit BitNet design, we achieve a representation equivalent to 1.58 bits. \text{BitNet b1.58} preserves all advantageous characteristics of the original 1-bit BitNet, such as its computation regime that minimizes reliance on multiplication in matrix operations, thus permitting aggressive computational optimizations. The energy usage remains comparable to that of the 1-bit BitNet, while \text{BitNet b1.58} surpasses FP16 LLMs in memory efficiency, throughput, and latency. In addition, \text{BitNet b1.58} offers two notable enhancements: First, by explicitly incorporating 0 into the set of weights, it can perform feature filtering, thus bolstering the model’s expressive capacity and yielding tangible performance improvements over purely 1-bit LLMs. Second, empirical evidence suggests that \text{BitNet b1.58} attains parity with full-precision (FP16) baselines in both perplexity and downstream evaluation tasks, beginning at the 3B parameter scale and under matching conditions such as model size and number of training tokens.

\section{BitNet b1.58}

\text{BitNet b1.58} extends the BitNet framework, a variant of the Transformer architecture that substitutes the standard \emph{nn.Linear} layer with the \emph{BitLinear} layer. This model is initialized from the beginning with weights quantized to 1.58 bits and activations quantized to 8 bits. While maintaining the foundational BitNet design, several alterations have been introduced, which are detailed below.

\paragraph{Quantization Function.} We restrict the weights to the set \{-1, 0, +1\} by applying an \emph{absmean} quantization scheme. This method begins by normalizing the weight matrix with its mean absolute value, after which each element is discretized to its closest integer within \{-1, 0, +1\}:

\begin{equation}
    \widetilde{W} = \mathrm{RoundClip}(\frac{W}{\gamma + \epsilon}, -1, 1),
\end{equation}

\begin{equation}
    \mathrm{RoundClip}(x, a, b) = \max(a, \min(b, \mathrm{round}(x))),
\end{equation}

\begin{equation}
    \gamma = \frac{1}{nm}\sum_{ij} |W_{ij}|.
\end{equation}

The quantization applied to activations mirrors the approach used in BitNet, with the exception that we forego rescaling the activations prior to non-linearity to the interval $[0, Q_b]$. Instead, activations for each token are rescaled within $[-Q_b, Q_b]$, eliminating the use of zero-point quantization. This adjustment simplifies both implementation and system-level optimizations, while our experimental results indicate that it produces negligible impact on performance.

\paragraph{LLaMA-alike Components.}

The LLaMA~\cite{llama, llama2} architecture has emerged as the standard template for open-source LLMs. To support and align with the open-source ecosystem, the architecture of \text{BitNet b1.58} incorporates components inspired by LLaMA. Notably, it employs RMSNorm~\cite{rmsnorm}, SwiGLU~\cite{swiglu}, rotary embedding~\cite{rope}, and eliminates all bias terms. These choices ensure that \text{BitNet b1.58} can be smoothly integrated with widely used open-source platforms, including Huggingface, vLLM~\cite{vllm}, and llama.cpp\footnote{\url{https://github.com/ggerganov/llama.cpp}}, with minimal overhead.

\section{Results}

We evaluated \text{BitNet b1.58} alongside our FP16 \text{LLaMA LLM} reimplementation at several parameter scales. To maintain evaluation consistency, both sets of models were pre-trained on the RedPajama dataset~\cite{redpajama} for a total of 100 billion tokens. For assessment, we used a diverse suite of zero-shot language understanding benchmarks, such as ARC-Easy~\cite{arc}, ARC-Challenge~\cite{arc}, Hellaswag~\cite{hellaswag}, Winogrande~\cite{winoGrande}, PIQA~\cite{piqa}, OpenbookQA~\cite{openbookqa}, and BoolQ~\cite{boolq}. Additionally, we provide validation perplexity statistics on the WikiText2~\cite{wikitext2} and C4~\cite{c4} corpora.

For system-level efficiency, we measured and compared GPU memory usage and inference latency between \text{LLaMA LLM} and \text{BitNet b1.58}. These metrics were collected using the FasterTransformer\footnote{\url{https://github.com/NVIDIA/FasterTransformer}} framework, known for its efficient LLM inference optimizations on GPUs. In particular, inference with \text{BitNet b1.58} benefited from integrating Ladder’s~\cite{ladder} 2-bit kernel. We reported the latency per generated token, as this directly impacts real-world deployment costs.

Perplexity outcomes and compute resource usage for both \text{BitNet b1.58} and \text{LLaMA LLM} are shown in Table~\ref{tab:ppl}. The findings indicate that \text{BitNet b1.58} achieves comparable perplexity to the full-precision \text{LLaMA LLM} at the 3B scale, with a speedup of 2.71$\times$ and 3.55$\times$ lower GPU memory requirements. More specifically, the 3.9B version of \text{BitNet b1.58} delivers 2.4$\times$ greater speed, reduces memory footprint by 3.32$\times$, and surpasses the 3B \text{LLaMA LLM} in performance.

A breakdown of zero-shot task accuracy is provided in Table~\ref{tab:zero-shot}. Our evaluation protocol followed the \emph{lm-evaluation-harness}\footnote{\url{https://github.com/EleutherAI/lm-evaluation-harness}} framework. The results demonstrate a shrinking performance difference between \text{BitNet b1.58} and \text{LLaMA LLM} as model size increases. Notably, \text{BitNet b1.58} matches the full-precision baseline’s results beginning at the 3B parameter setting. In line with the perplexity metrics, end-task evaluations confirm that the 3.9B variant of \text{BitNet b1.58} not only uses less memory and runs faster but also exceeds the performance of the 3B \text{LLaMA LLM}. This highlights \text{BitNet b1.58} as a clear Pareto improvement compared to leading LLM baselines.

\paragraph{Memory and Latency}

The model was further scaled to 7B, 13B, and 70B parameters, and we assessed the corresponding computational cost. Figure~\ref{fig:latency-memory-energy} presents how latency and memory usage change with model size, indicating a greater acceleration in speed as models become larger. Specifically, \text{BitNet b1.58} at 70B achieves a speed-up of 4.1$\times$ over the \text{LLaMA LLM} reference. This improvement stems from the increasing computational burden of \emph{nn.Linear} with scaling model size. The observed memory usage aligns with this pattern, since the full-precision embedding's memory footprint diminishes proportionally in larger architectures. All latency and memory metrics were gathered using a 2-bit kernel implementation, suggesting further potential for future cost reduction through kernel optimization.

\paragraph{Energy}

We also evaluated the energy consumption resulting from arithmetic operations in both \text{BitNet b1.58} and \text{LLaMA LLM}. Matrix multiplication, being the predominant factor in the total computation load for LLMs, was the main focus of our analysis. As depicted in Figure~\ref{fig:energy}, energy usage is chiefly attributed to INT8 additions in \text{BitNet b1.58}, in contrast to the FP16 additions and multiplications comprising \text{LLaMA LLM}. Drawing on the energy estimation models from~\cite{energycost, pokebnn}, it was found that \text{BitNet b1.58} offers a 71.4$\times$ reduction in matrix multiplication energy consumption for 7nm semiconductor processes. We also reported full end-to-end energy requirements for sequences of 512 tokens, demonstrating that as the model grows larger, \text{BitNet b1.58} yields even greater energy efficiency relative to the FP16 \text{LLaMA LLM} baseline. This enhanced efficiency arises because the share of computational cost attributed to \emph{nn.Linear} expands with increased model size, while other components contribute less to the overall energy usage in larger networks.

\paragraph{Throughput}

We evaluate and contrast the throughput performance of \text{BitNet b1.58} and \text{LLaMA LLM}, each featuring 70B parameters, by utilizing two 80GB A100 GPUs and adopting pipeline parallelism~\cite{gpipe}. This setup ensures that \text{LLaMA LLM} 70B can operate within the hardware constraints. The batch size was incrementally increased until the GPUs reached their maximum memory capacity, with the sequence length fixed at 512. According to Table~\ref{tab:throughput}, the \text{BitNet b1.58} 70B model accommodates a batch size up to 11 times greater than that of the \text{LLaMA LLM}, which yields an 8.9-fold improvement in throughput.

\textbf{\text{BitNet b1.58} introduces a novel scaling law relating model efficiency and inference resource requirements.} For further illustration, we present a set of size equivalences between 1.58-bit and 16-bit models as deduced from Figure~\ref{fig:latency-memory-energy} and \ref{fig:energy}.

\begin{itemize}
    \item The 13B \text{BitNet b1.58} surpasses the 3B FP16 LLM in latency, memory consumption, and energy usage.
    \item The 30B \text{BitNet b1.58} exhibits superior latency, memory efficiency, and power consumption compared to the 7B FP16 LLM.
    \item The 70B \text{BitNet b1.58} outperforms the 13B FP16 LLM regarding latency, memory requirements, and energy expenditure.
\end{itemize}

\paragraph{Training with 2T Tokens}

Token count during training plays a pivotal role for large language models. To assess the scalability of \text{BitNet b1.58} with respect to training tokens, we trained a \text{BitNet b1.58} model using 2T tokens, employing the data preparation methodology of StableLM-3B~\cite{StableLM-3B-4E1T}—a leading open-source 3B model. Both \text{BitNet b1.58} and StableLM-3B were evaluated on a benchmark comprising Winogrande~\cite{winoGrande}, PIQA~\cite{piqa}, SciQ~\cite{sciq}, LAMBADA~\cite{lambada}, and ARC-easy~\cite{arc}. Zero-shot accuracies are reported in Table~\ref{tab:2t}. Where tasks included both accuracy and normalized accuracy, we computed their mean. Results for StableLM 3B at 2T tokens were sourced directly from its technical documentation. Our experiments demonstrate that \text{BitNet b1.58} achieves higher scores on all evaluated tasks, signifying that 1.58-bit LLMs also exhibit strong generalization performance.

\section{Discussion and Future Work}

\textbf{1-bit Mixture-of-Experts (MoE) LLMs}

Mixture-of-Experts (MoE) architectures have emerged as an efficient solution to reduce computational complexity in LLMs, especially by lowering FLOPs. However, their large memory requirements and considerable inter-chip communication overhead present substantial barriers to widespread adoption. The introduction of 1.58-bit LLMs offers promising remedies to these issues. By minimizing the memory demands, it becomes feasible to decrease the number of hardware devices necessary for MoE deployment. Additionally, communication costs for transferring activations across networked devices are substantially reduced. Ideally, if these models can be entirely allocated to a single chip, such overhead could be eliminated.

\textbf{Native Support of Long Sequence in LLMs}

The ability to process extended sequences natively is increasingly vital for modern LLM applications. One of the primary obstacles in handling longer contexts during inference lies in the memory required to maintain the KV caches. \text{BitNet b1.58} advances this area by compressing activations from 16 bits to 8 bits, thereby enabling a twofold increase in sequence length under identical memory constraints. Further, for 1.58-bit LLMs, there exists potential to compress activation representations losslessly down to 4 bits or even less—a direction we reserve for subsequent research.

\textbf{LLMs on Edge and Mobile}

Implementing 1.58-bit LLMs could dramatically enhance language model performance on edge and mobile platforms, where constraints on memory and computational resources typically limit the feasibility of running large models. The substantially decreased memory and power consumption associated with 1.58-bit LLMs makes it possible to deploy these models on-device, unlocking a range of applications previously considered unachievable. This, in turn, extends the functional scope and innovation potential of edge and mobile computing environments. Furthermore, since 1.58-bit LLMs are inherently better suited to CPU execution—which remains the dominant processing architecture in mobile and edge settings—\text{BitNet b1.58} can deliver efficient and performant inference in these contexts.

\textbf{New Hardware for 1-bit LLMs}

Recent developments, such as those demonstrated by~Groq\footnote{\url{https://groq.com/}}, highlight significant advances and opportunities for crafting dedicated hardware solutions like LPUs tailored for LLM workloads. Building upon this progress, we propose the development of hardware and computing systems intentionally designed for 1-bit LLMs, taking full advantage of the novel computation patterns introduced by BitNet~\cite{bitnet}.